\documentclass[11pt]{article}

\usepackage[margin=1.1in]{geometry}
\usepackage{amsmath,amssymb,amsthm,mathrsfs}
\usepackage[colorlinks=true,linkcolor=black,urlcolor=black,citecolor=black,
  hypertexnames=false]{hyperref}

\newtheorem{theorem}{Theorem}
\newtheorem{lemma}[theorem]{Lemma}
\newtheorem{proposition}[theorem]{Proposition}
\newtheorem{corollary}[theorem]{Corollary}
\theoremstyle{definition}
\newtheorem{definition}[theorem]{Definition}
\theoremstyle{remark}
\newtheorem{remark}[theorem]{Remark}

\newcommand{\C}{\mathbf C}
\newcommand{\Q}{\mathbf Q}
\newcommand{\Z}{\mathbf Z}
\newcommand{\PP}{\mathbf P}

\title{An unconditional $m=1$ theorem via a rank-two atomic residue invariant}
\author{Tavis Rudd}
\date{August 2026}

\begin{document}
\maketitle

\begin{quote}
\emph{Editorial note (not part of the argument).} The text below is transcribed
verbatim from an external review by Sol dated 15 August 2026; the wording,
section headings, equation numbering and boxed statements are the reviewer's,
and nothing has been paraphrased, corrected or added. The argument is
\textbf{unverified} as of transcription, and a separate verification report is
being prepared; until that report exists, no claim here should be relied on.
The argument imports the ordinary, non-enhanced Hodge-atom results of
Katzarkov--Kontsevich--Pantev--Yu as external inputs, and those results are
presently available only as an arXiv preprint.
\end{quote}

\bigskip

Yes. Subject to the ordinary Hodge-atom theorems of
Katzarkov--Kontsevich--Pantev--Yu as external inputs, the following gives an
unconditional proof of the $m=1$ theorem. It does not use Hypothesis 4.7H, the
pro-Laurent gauge, divisor tagging, or transport of framed monodromy through
weak factorization.

I will distinguish the imported KKP(Y) results from the new argument.

\subsection*{Theorem}

Let $X\subset\PP^4_\C$ be any smooth cubic threefold. Then $X\times\PP^1$ is
irrational.

The only external atom-theoretic inputs are the ordinary, non-enhanced
Hodge-atom results of KKP(Y): spectral decomposition of the non-archimedean
$A$-model $F$-bundle; the map from Hodge atoms to geometric atomic $F$-bundles;
invariance of the underlying Hodge representation; the projective-bundle
chemical formula; and the non-rationality criterion. We do not use their
enhanced Serre-automorphism argument in Example 6.21.

\section{A rank-two invariant of an ordinary geometric Hodge atom}

KKP(Y)'s $F$-bundles have, in their loop coordinate $u$, poles satisfying
$\nabla_{u^2\partial_u}$, $\nabla_{u\xi}$ regular for every base vector field
$\xi$. Their $A$-model connection is
\[
  \nabla_{u\partial_u}=u\partial_u-u^{-1}\mathrm{Eu}\star+\mathrm{Gr},
  \qquad
  \nabla_\xi=\xi+u^{-1}C_\xi,
\]
with the obvious form for Novikov directions.

Let $\alpha$ be a geometric Hodge atom and suppose its even part $E$ has rank
two. On the Hodge-fixed base its Euler operator has a single eigenvalue by the
definition of an atomic $F$-bundle.

Write $U$ for the leading Euler operator on $E|_{u=0}$, and put
$\lambda=\tfrac12\operatorname{tr}U$, $N=U-\lambda I$.

Since $U$ has only the eigenvalue $\lambda$ and has rank two,
\begin{equation}
  N^2=0.
  \tag{1}
\end{equation}

We shall treat the case in which $N$ is not identically zero on the connected
spectral component representing $\alpha$.

Center the exponential factor. Replace the connection by
\begin{equation}
  \nabla^\circ=\nabla-d(\lambda/u)\,I.
  \tag{2}
\end{equation}

This is only a device for defining the invariant. It uses the same underlying
lattice and removes the scalar exponential factor.

In a local frame, horizontal sections of the even factor satisfy
\begin{equation}
  u\partial_u y=A(u)\,y,\qquad \delta y=B_\delta(u)\,y,
  \tag{3}
\end{equation}
where, because of the $F$-bundle pole conditions,
\begin{equation}
  A(u)=u^{-1}N+A_0+uA_1+u^2A_2+\cdots,
  \tag{4}
\end{equation}
\begin{equation}
  B_\delta(u)=u^{-1}C_\delta+C_{\delta,0}+uC_{\delta,1}+\cdots.
  \tag{5}
\end{equation}

Flatness gives
\begin{equation}
  \delta A-u\partial_uB_\delta+[A,B_\delta]=0.
  \tag{6}
\end{equation}

The Poincar\'e pairing forces $A_0$ to preserve $\ker N$.

The Poincar\'e pairing is Frobenius for quantum multiplication;
$\mathrm{Eu}\star$ is self-adjoint and $\mathrm{Gr}$ anti-self-adjoint. Hence it
is horizontal for the quantum connection. KKP(Y) record exactly this
calculation; it is algebraic in the quantum product and therefore applies
equally to their formal/non-archimedean $A$-model.

The spectral factors inherit nondegenerate pairings; more abstractly, KKP(Y)
note that pairings are compatible with spectral decomposition. On the even
rank-two factor write its matrix as $P(u)=P_0+uP_1+u^2P_2+\cdots$.

Here $P_0$ is nondegenerate and symmetric. Centering by (2) does not alter
horizontality, because the scalar terms at $u$ and $-u$ cancel.

Horizontality in the $u$-direction says
\begin{equation}
  u\partial_uP(u)+A(u)^TP(u)+P(u)A(-u)=0.
  \tag{7}
\end{equation}

The coefficient of $u^{-1}$ is
\begin{equation}
  N^TP_0=P_0N.
  \tag{8}
\end{equation}

Thus $N$ is self-adjoint for $P_0$.

On the locus where $N\ne0$, (1) implies $L:=\operatorname{im}N=\ker N$ is a
line. Moreover it is isotropic: if $v=Nx$, $w=Ny$, then
$(v,w)=(Nx,Ny)=(x,N^2y)=0$.

As $P_0$ is nondegenerate on a two-dimensional space,
\begin{equation}
  L^\perp=L.
  \tag{9}
\end{equation}

The coefficient of $u^0$ in (7) is
\begin{equation}
  A_0^TP_0+P_0A_0+N^TP_1-P_1N=0.
  \tag{10}
\end{equation}

Take $0\ne x$ in $L=\ker N$ and sandwich (10) between $x$'s. The terms involving
$N$ vanish, giving $2(A_0x,x)=0$.

Therefore $A_0x\in L^\perp=L$. Hence
\begin{equation}
  \boxed{\,A_0L\subset L.\,}
  \tag{11}
\end{equation}

This is the crucial pairing consequence. Notice that no pairing-preserving
block-splitting gauge has been assumed.

\section{The canonical elementary modification}

Define, on the open locus $N\ne0$,
\begin{equation}
  E^\sharp=\{\,s\in E:\ s\bmod u\in L\,\}.
  \tag{12}
\end{equation}

This is intrinsic: it is the elementary modification of the original $F$-bundle
lattice along the canonical line $L=\operatorname{im}N$.

Choose locally a basis $e_1,e_2$ such that
\begin{equation}
  L=\langle e_1\rangle,\qquad Ne_1=0,\qquad Ne_2=\nu e_1,\qquad
  \nu\in\mathcal O^\times.
  \tag{13}
\end{equation}

Then $E^\sharp$ has basis $e_1,ue_2$. Thus the change-of-lattice matrix is
$S=\operatorname{diag}(1,u)$.

The $u$-matrix in this lattice is
\begin{equation}
  A^\sharp=S^{-1}AS-S^{-1}u\partial_uS.
  \tag{14}
\end{equation}

Since $S^{-1}E_{12}S=uE_{12}$, $S^{-1}E_{21}S=u^{-1}E_{21}$, the leading term
$u^{-1}N=\nu u^{-1}E_{12}$ becomes the regular term $\nu E_{12}$.

The only possible pole coming from $A_0$ would be its $E_{21}$-entry. But (11)
says precisely that this entry is zero. Every term $u^mA_m$, $m\ge1$, can lose
at most one power of $u$, and hence remains regular.

Consequently
\begin{equation}
  A^\sharp(u)=R+uR_1+u^2R_2+\cdots.
  \tag{15}
\end{equation}

Thus the elementary-modified $u$-connection is regular singular --- in fact
$u\partial_u$ has a regular matrix --- and its residue matrix
\begin{equation}
  R=A^\sharp(0)
  \tag{16}
\end{equation}
is intrinsically defined up to conjugacy.

Indeed, under a regular change of frame $F(u)$ in $\mathrm{GL}_2(\mathcal
O[[u]])$, $R\to F(0)^{-1}RF(0)$, because $u\partial_uF=0\bmod u$.

So
\begin{equation}
  \chi^\sharp_\alpha(T):=\det(TI-R)
  \tag{17}
\end{equation}
is intrinsic wherever $N\ne0$.

\section{Flatness makes $\chi^\sharp_\alpha$ constant}

We now prove the point that replaces the missing ``monodromy is constant on
connected spectral components.''

From (6), the $u^{-2}$-coefficient gives
\begin{equation}
  [N,C_\delta]=0.
  \tag{18}
\end{equation}

The $u^{-1}$-coefficient is
\begin{equation}
  \delta N+C_\delta+[N,C_{\delta,0}]+[A_0,C_\delta]=0.
  \tag{19}
\end{equation}

Taking traces and using $\operatorname{tr}N=0$ gives
\begin{equation}
  \operatorname{tr}C_\delta=0.
  \tag{20}
\end{equation}

For a nonzero rank-one nilpotent $2\times2$ matrix, its commutant is
$\mathcal OI\oplus\mathcal ON$. Hence (18) and (20) imply
\begin{equation}
  C_\delta=q_\delta N
  \tag{21}
\end{equation}
for some local scalar $q_\delta$.

Transform the base equation to the modified lattice. Because $u^{-1}q_\delta N$
becomes regular, the only possible pole is produced by the $E_{21}$-entry of
$C_{\delta,0}$. Therefore
\begin{equation}
  B^\sharp_\delta=u^{-1}K_\delta+G_\delta+O(u),\qquad
  K_\delta=k_\delta E_{21}.
  \tag{22}
\end{equation}

Flatness remains
\begin{equation}
  \delta A^\sharp-u\partial_uB^\sharp_\delta+[A^\sharp,B^\sharp_\delta]=0.
  \tag{23}
\end{equation}

Its $u^{-1}$-coefficient is
\begin{equation}
  K_\delta+[R,K_\delta]=0.
  \tag{24}
\end{equation}

Because the constant $E_{12}$-entry of $R$ comes from $N$, it is the unit $\nu$.
Thus
\begin{equation}
  R=\begin{pmatrix}a&\nu\\c&d\end{pmatrix}.
  \tag{25}
\end{equation}

For $K_\delta=k_\delta E_{21}$,
$[R,K_\delta]=k_\delta\begin{pmatrix}\nu&0\\d-a&-\nu\end{pmatrix}$.

The diagonal entries of the left side of (24) are therefore $\nu k_\delta$,
$-\nu k_\delta$. Since $\nu$ is a unit,
\begin{equation}
  \boxed{\,K_\delta=0.\,}
  \tag{26}
\end{equation}

So after the canonical elementary modification the connection is regular not
only in $u$ but also in every base direction.

The coefficient of $u^0$ in (23) now gives $\delta R+[R,G_\delta]=0$, or
\begin{equation}
  \boxed{\,\delta R=[G_\delta,R].\,}
  \tag{27}
\end{equation}

It follows immediately that $\delta\operatorname{tr}R=0$,
$\delta\operatorname{tr}(R^2)=0$, and therefore $\delta\det R=0$.

Hence
\begin{equation}
  \boxed{\,\delta\chi^\sharp_\alpha(T)=0\,}
  \tag{28}
\end{equation}
for every tangent vector $\delta$.

This is finite-dimensional isomonodromy. No formal bulk gauge has appeared.

\section{Crossing points where $N=0$}

This is the point at which the memo's atom route stopped.

Let $\mathfrak U$ be the connected component of the unramified spectral cover
representing the atom. KKP(Y) define local atoms precisely as connected
components of this cover and identify atomic $F$-bundle germs at two points when
their spectral lifts lie in the same connected component.

On the open locus $\mathfrak U^\times=\{N\ne0\}$ the preceding construction is
canonical.

If $N$ is not generically zero, the line $L=\operatorname{im}N$ is defined over
the meromorphic function field of $\mathfrak U$. Thus $E^\sharp$, $R$, and the
two coefficients of $\chi^\sharp_\alpha$ are meromorphic objects on
$\mathfrak U$.

On $\mathfrak U^\times$, (28) says those meromorphic functions are locally
constant.

The reduced spectral cover over the locus of maximal number of distinct
eigenvalues is locally a disjoint union of eigenvalue branches; consequently a
connected component is smooth and locally irreducible. Hence a connected
component is integral. This is also exactly the setting in which KKP(Y)
analytically continue one atomic factor along a connected spectral component.

Choose any nonempty open subset of $\mathfrak U^\times$. There
$\chi^\sharp_\alpha(T)=T^2+c_1T+c_0$ with $c_0,c_1$ constants. Since the two
coefficients are meromorphic functions on the integral space $\mathfrak U$,
equality with constants on a nonempty open forces equality as meromorphic
functions everywhere.

Therefore distinct components of $\mathfrak U^\times$, even if separated by
points at which $N=0$, cannot acquire different values.

So:
\begin{equation}
  \boxed{\ \chi^\sharp_\alpha\ \text{is an invariant of the entire connected
  spectral component.}\ }
  \tag{29}
\end{equation}

No assertion that $N$ itself remains nonzero pointwise is needed.

If $N$ is generically zero, record that as a separate type. We will only use
(17) for atoms for which $N$ is generically nonzero.

\subsection*{Invariance under atomic $F$-bundle isomorphism}

Let $F:(E,\nabla)\to(E',\nabla')$ be an isomorphism of geometric atomic
$F$-bundles. Because $F$ is regular and invertible at $u=0$,
$N'=F(0)^{-1}NF(0)$, so it carries $L=\operatorname{im}N$ to
$L'=\operatorname{im}N'$, and hence carries $E^\sharp$ to $E'^\sharp$.

The induced residue matrices are conjugate. Therefore
\begin{equation}
  \chi^\sharp_{\alpha'}=\chi^\sharp_\alpha.
  \tag{30}
\end{equation}

Together, (29) and (30) show that $\chi^\sharp$ descends through both
equivalences in KKP(Y)'s Definition 5.21.

Finally, Proposition 5.22 maps every abstract Hodge atom to a well-defined
geometric atomic $F$-bundle equivalence class. Thus:

\medskip
\noindent\textbf{Rank-two atomic residue lemma.} For a Hodge atom whose even
geometric atomic $F$-bundle has rank two and whose centered leading operator is
generically nonzero nilpotent, the characteristic polynomial
$\chi^\sharp_\alpha(T)$ of the elementary-modified residue is a well-defined
invariant of the ordinary Hodge atom.
\medskip

This uses the Poincar\'e pairing only to prove (11). The pairing is not being
added to the atom as enhancement data.

That distinction is important.

\section{The cubic atom}

KKP(Y), in Example 6.21, state for every smooth cubic threefold that at the
hyperplane point the atomic composition consists of two one-dimensional
nonzero-eigenvalue atoms and one zero-eigenvalue atom $\alpha(X)$, with no
further local spectral splitting.

Its even rank is two. For completeness, here is the dimension calculation.

For a smooth cubic threefold, $c(TX)=(1+H)^5/(1+3H)=1+2H+4H^2-2H^3$, so
$\chi_{\mathrm{top}}(X)=\int_Xc_3(TX)=-2\int_XH^3=-6$.

By Lefschetz and Poincar\'e duality, $b_0=b_2=b_4=b_6=1$, $b_1=b_5=0$.

Thus
\begin{equation}
  4-b_3=-6,\qquad b_3=10.
  \tag{31}
\end{equation}

Hence total even and odd dimensions are respectively $4$, $10$.

The two other cubic atoms are one-dimensional. Every Hodge atom has a nonzero
even unit vector, so a one-dimensional atom is necessarily even. Consequently
\begin{equation}
  \operatorname{rk}\alpha(X)_{\bar0}=2,\qquad
  \operatorname{rk}\alpha(X)_{\bar1}=10.
  \tag{32}
\end{equation}

These parity dimensions themselves are invariants of the atom by KKP(Y)'s
Proposition 5.23.

\subsection*{Compute $\chi^\sharp_{\alpha(X)}$}

The paper's unconditional cubic calculation block-diagonalizes the original-$z$
small even connection. On the zero generalized eigenspace it obtains
\[
  z^2\partial_zY=(J_0+zD_0+z^2E_0+O(z^3))\,Y,
\]
with $J_0=\begin{pmatrix}0&2\\0&0\end{pmatrix}$,
$D_0=\operatorname{diag}(-19/18,19/18)$, and $(E_0)_{21}=-8/81$.

This part of Proposition 4.12 is unconditional.

Dividing by $z$, $z\partial_zY=(z^{-1}J_0+D_0+zE_0+\cdots)Y$.

Here $\lambda=0$ and $N=J_0=2E_{12}$, $L=\langle e_1\rangle$.

The canonical elementary modification therefore uses
$S=\operatorname{diag}(1,z)$.

Its residue is
\begin{equation}
  R_X=2E_{12}+D_0-\tfrac{8}{81}E_{21}-\operatorname{diag}(0,1)
  =\begin{pmatrix}-19/18&2\\-8/81&1/18\end{pmatrix}.
  \tag{33}
\end{equation}

Thus $\operatorname{tr}R_X=-1$, $\det R_X=5/36$, and hence
\begin{equation}
  \boxed{\ \chi^\sharp_{\alpha(X)}(T)=T^2+T+\tfrac{5}{36}
  =\bigl(T+\tfrac16\bigr)\bigl(T+\tfrac56\bigr).\ }
  \tag{34}
\end{equation}

This is the same $-1/6$, $-5/6$ pair occurring in the paper's framed-monodromy
calculation, but here it is being used as the residue spectrum of a canonical
modification of the intrinsic atomic lattice, not as monodromy transported
through a bulk gauge.

\section{The only possible curve competitor is genus five}

Suppose $\alpha(X)$ were represented by a smooth projective curve $C$.

For $C=\PP^1$, $\PP^1=\PP_{\mathrm{pt}}(\mathcal O^{\oplus2})$, so the
projective-bundle chemical formula says its atoms are two point atoms. Thus
$\alpha(X)$ cannot occur.

Let $g(C)\ge1$. Then $K_C$ is nef, so KKP(Y)'s Lemma 5.24 says $C$ has a single
Hodge atom. Its underlying representation is therefore all of $H^\bullet(C)$.

Its parity dimensions are $(2,2g)$. Comparison with (32) forces
\begin{equation}
  g=5.
  \tag{35}
\end{equation}

This recovers directly the observation in KKP(Y) that the cubic atom ``looks
like'' the genus-five curve atom at the ordinary Hodge-representation level.

We now distinguish them using $\chi^\sharp$.

\subsection*{Compute the curve residue invariant}

Let $p\in H^2(C)$ satisfy $\int_Cp=1$. For $g\ge1$, there are no nonconstant
genus-zero maps to $C$, so its genus-zero quantum product is the classical cup
product.

Put $a=2-2g$. Then $c_1(TC)=ap$.

In the ordered even basis $1,p$, $c_1(TC)\cup(-)=aE_{21}$.

For a curve, $\mathrm{Gr}=\operatorname{diag}(-1/2,1/2)$.

Using the $A$-model connection formula of KKP(Y), horizontal sections satisfy
\begin{equation}
  u\partial_uY=\bigl[\,u^{-1}aE_{21}+\operatorname{diag}(1/2,-1/2)\,\bigr]Y.
  \tag{36}
\end{equation}

This agrees with Cai's direct curve computation.

For $g>1$, $a\ne0$, and $N=aE_{21}$, $L=\langle e_2\rangle$.

The elementary modification has basis $ue_1,e_2$, so now
$S=\operatorname{diag}(u,1)$.

The modified residue is
\begin{equation}
  R_C=aE_{21}+\operatorname{diag}(1/2,-1/2)-\operatorname{diag}(1,0)
  =aE_{21}-\tfrac12I.
  \tag{37}
\end{equation}

Therefore
\begin{equation}
  \boxed{\ \chi^\sharp_C(T)=\bigl(T+\tfrac12\bigr)^2.\ }
  \tag{38}
\end{equation}

For a genus-five curve in particular,
\begin{equation}
  \bigl(T+\tfrac12\bigr)^2\ne T^2+T+\tfrac{5}{36}.
  \tag{39}
\end{equation}

Hence
\begin{equation}
  \boxed{\ \alpha(X)\ \text{cannot be represented by a curve.}\ }
  \tag{40}
\end{equation}

\section{It cannot be represented by a surface}

Suppose $\alpha(X)$ occurs in the atomic composition of a smooth projective
surface $S$.

Blowing down $(-1)$-curves to a minimal model only removes point atoms: for a
point blowup the ordinary atom formula is
\begin{equation}
  \mathrm{CF}(\mathrm{Bl}_pS)=\mathrm{CF}(S)+\mathrm{CF}(p).
  \tag{41}
\end{equation}

Since $\alpha(X)$ has total rank $12$, it is not a point atom. Therefore it must
already occur on a minimal model.

If the minimal surface has nef canonical bundle, Lemma 5.24 says its atomic
composition consists of a single atom, whose underlying representation is its
entire cohomology. Its even part has dimension at least three: $H^0(S)$,
$\langle H\rangle\subset H^2(S)$, $H^4(S)$ give three independent even classes
for any ample divisor $H$.

But $\alpha(X)$ has even rank two by (32). Contradiction.

If the minimal surface does not have nef canonical class, the surface
classification gives either $\PP^2$ or a geometrically ruled surface
$\PP_C(V)$. KKP(Y) use precisely this low-dimensional surface classification in
their atom analysis.

For $\PP^2$, $\mathrm{CF}(\PP^2)=3\,\mathrm{CF}(\mathrm{pt})$.

For a ruled surface, $\mathrm{CF}(\PP_C(V))=2\,\mathrm{CF}(C)$.

These are immediate applications of the projective-bundle formula. Thus their
atoms are point or curve atoms, already excluded.

Consequently
\begin{equation}
  \boxed{\ \alpha(X)\notin\mathrm{HAtoms}_{\dim\le2}.\ }
  \tag{42}
\end{equation}

\section{One stabilization}

Now $X\times\PP^1=\PP_X(\mathcal O_X^{\oplus2})$.

KKP(Y)'s projective-bundle formula gives
\begin{equation}
  \mathrm{CF}(X\times\PP^1)=2\,\mathrm{CF}(X).
  \tag{43}
\end{equation}

Hence the atom $\alpha(X)$ occurs, with positive multiplicity, in the atomic
composition of $X\times\PP^1$.

The latter is a smooth projective fourfold. KKP(Y)'s Proposition 5.30 says that
if a smooth $d$-fold contains a Hodge atom not realizable by a smooth projective
variety of dimension at most $d-2$, the variety is not rational.

Here $d=4$, and (42) says precisely that
$\alpha(X)\notin\mathrm{HAtoms}_{\dim\le2}$.

Therefore
\[
  \boxed{\ X\times\PP^1\ \text{is irrational.}\ }
\]

This proves the theorem.

\section*{What this means for \texttt{test3.pdf}}
\addcontentsline{toc}{section}{What this means for \texttt{test3.pdf}}

This gives a replacement for the current conditional proof of Theorem 1.1:
\[
  \boxed{\ \text{Theorem 1.1 can be made unconditional without proving 4.7H.}\ }
\]

The new ingredient is the rank-two atomic residue lemma above. The cubic
calculation already present in Proposition 4.12 supplies its value
$T^2+T+5/36$.

The competing genus-five curve has value $(T+1/2)^2$.

The ordinary KKP(Y) atom ledger then handles the stabilization and rationality
obstruction.

Hypothesis 4.7H would still matter if you want the much stronger
summand-by-summand operation formulas of Proposition 4.7, the general $\nu_6$
birational-invariance theorem, or the current genus-eight argument. But it is no
longer needed for the cubic $m=1$ headline theorem.

There is one source-level caveat rather than a mathematical hypothesis: KKP(Y)
v2 is presently an arXiv preprint. This proof imports its ordinary atom
decomposition/projective-bundle/non-rationality machinery as theorem. It does
not import the enhanced Serre/integral-structure machinery used in their own
Example 6.21.

\end{document}
